% page 570 du fac-simile (image p-569 du scan)
% bloc 160.0 x 250.6 mm ; marge gauche 8.0 mm
\pgc{-12.700mm}{0.000mm}{
\l{\cel{3}562}
\l{}
\l{\sou{suolo}. I. (\sou{anat.}) Suba parto di la hufo, che ula animali}
\l{\cel{3}(kavalo, asno, mulo, cervo, e c.), plako korna inter la}
\l{\cel{3}bordumo infra di la parieto, la forketo e la apogarki qui}
\l{\cel{3}cirkondas lu. - II. Parto di la pedovesto, qua esas sub}
\l{\cel{3}la plando, por protektar ica. - DEIS.}
\l{}
\l{\sou{supo}.\cel{2}Buliono quan onu varsas sur pano-lonchi dina, disho}
\l{\cel{3}per qua komencas, ordinare, la repasto. - DEFIRS.}
\l{}
\l{\sou{supear}. (\sou{netrans}.)Repastar ye kloko (generale) kelke pos nok-}
\l{\cel{3}to-mezo, kande onu retrovenas de festo, de balo, e c.- DEF}
\l{}
\l{\sou{super}.\cel{2}Prepoziciono : Qua esas plu alta kam la suprajo e sen}
\l{\cel{3}kontakto kun ica. -\cel{1}\sou{anke metaf}.\cel{4}EFIS.}
\l{}
\l{\sou{superba}. Qua estimas su tante ke lu situas su super sua kun-}
\l{\cel{3}homi ed, pro lo, exajeras sua revendiki sociala, mondumala.}
\l{\cel{3}- FIS.}
\l{}
\l{\sou{superfetaciono}. (\sou{fiziol.}) Formaceso di feto duesma super feto}
\l{\cel{3}unesma, konceptita antee. - DEFIS.}
\l{}
\l{\sou{superflua}. Qua esas ultre lo necesa, exter lo utila. - EFIS.}
\l{}
\l{\sou{superioro}. Persono qua esas super altra en hierarkio, e qua ha-}
\l{\cel{3}vas la yuro obediigar su da ta altra. - DEFIS.}
\l{}
\l{\sou{superlativo}. (\sou{gram.}) La grado maxim alta di la eso-maniero quan}
\l{\cel{3}expresas epiteto. - DEFIS.}
\l{}
\l{\sou{supernumerario}. Employato quan administranti, komerco-firmo,}
\l{\cel{3}engajas suplemente, varte di la vako di ofico che su. DEFIS.}
\l{}
\l{\sou{superstico}.\cel{2}I. Kredo, praktiko a qua religio miskomprenata}
\l{\cel{3}donas karaktero sakra. - II. Eceso di skrupulo pri la egar-}
\l{\cel{3}do di ula kozi. - EFIS.}
\l{}
\l{\sou{supino}. (\sou{gram. Latina}) Tempo di infinitivo qua esas analoga,}
\l{\cel{3}pri la formo, a participo pasintala, e, pri la senco, a sub-}
\l{\cel{3}stantivo obtenita de la radiko di verbo. - DEFIS.}
\l{}
\l{\sou{suplantar}. (\sou{trans.})\cel{2}Eviktar ulu e substitucar su a lu. -EFIS.}
\l{}
\l{\sou{suplear}. (\sou{trans.)}\cel{3}Remplasar fristoze, pleante la sama rolo,}
\l{\cel{3}(to quo ne povas, ta qua ne povas, plear sua rolo).--EFIS.}
\l{}
\l{\sou{suplemento}. To quo donesas o pagesas ultre lo debata. - DEFIS.}
\l{}
\l{\sou{suplikar}. (\sou{trans.}) (\sou{pri, por}).Pregar ye maniero tre humila e}
\l{\cel{3}tre insistoza, ye la nomo di ulo kara e sakra. - EFIS.}
\l{}
\l{\sou{suportar}. (\sou{trans.}) I. (a) Havar sur su (kozo di qua onu portas}
\l{\cel{3}la tota pezo) talo ke onu impedas olua falo ; (b) (\sou{metaf}.)}
\l{\cel{3}Havar sur su la tota kargajo, la tota embaraso di (ulo).-}
\l{\cel{3}II. (a) Recevar sur su sen\cel{2}shancelar de lo (la tota esforco}
\l{\cel{50}...}
}
